% !TeX root = ../../Book4.tex
% 纲要标题：### 12.4 分类的唯一性
% 纲要位置：几何学计划纲要.md，第 3947 行。
% 正文状态：初稿；尚未正式审核。

% 纲要内容开始
% 
% * 局部标准模型的比较
% * 自由轨道部分的粘合
% * 边界轨道的延拓
% * 等变辛同构与平移自由
% 
% 纲要内容结束

\section{分类的唯一性}
\label{g4:sec:1204}

由多面体作约化得到一个流形，还没有说明给定的流形一定与它相同。唯一性需要比较两套局部模型，并把局部同构粘成整体同构。多面体内部的纤维是自由环面轨道；边界上有圆方向收缩，角坐标会失效。因此我们先确定全部稳定子和边界模型，再把粘合误差写成底空间上的光滑 Hamilton 函数。凸多面体的可缩性将消去这些误差。

\subsection{局部标准模型的比较}
\label{g4:subsec:1204:01}
本节固定 \(T=\mathfrak t/(2\pi\Lambda)\)，其中 \(\Lambda\cong\mathbb Z^n\)。辛环面流形是带有效 Hamilton \(T\) 作用的紧连通辛流形 \((M^{2n},\omega)\)，并指定矩映射 \(\mu:M\to\mathfrak t^*\)。先由一般轨道标准形推出半维作用的特殊性质。

\begin{proposition}\label{g4:prop:1204:toric-orbits}
	设 \((M^{2n},\omega,T,\mu)\) 为辛环面流形。则每个稳定子 \(H=T_x\) 都连通。若 \(\dim H=h\)，辛切片的复维数为 \(h\)，其 \(h\) 个权构成 \(H\) 的特征格的一组整数基。每条矩纤维恰是一条 \(T\) 轨道，矩映射诱导同胚
	\[
		M/T\xrightarrow{\ \cong\ }P=\mu(M).
	\]
	\(P\) 是满维 Delzant 多面体；余维 \(h\) 面的相对内部对应稳定子维数为 \(h\) 的轨道。
\end{proposition}
\begin{proof}
	\Cref{g4:thm:1201:orbit-normal-form}给出 \(x\) 的轨道附近的模型
	\[
		T\times_H(\mathfrak h^\circ\times V),\qquad
		\dim_{\mathbb R}V=2n-2\dim(T/H)=2h.
	\]
	给 \(V\) 取 \(H\) 不变的相容复结构。紧交换群的酉表示可同时对角化，因而有 \(h\) 个字符 \(\chi_j:H\to S^1\)。它们给出的同态 \(\chi:H\to(S^1)^h\) 是单射。事实上，核中的元素在 \(V\) 上恒等，因 \(T\) 交换，它在模型的轨道方向及 \(\mathfrak h^\circ\) 上也恒等，所以在 \(M\) 的一个开集上恒等。

	为说明这迫使该元素全局恒等，平均任一 Riemann 度量，使 \(T\) 作用等距。第三卷定理2.1.2的测地初值唯一性使等距映射与指数映射交换，命题2.1.4又说明指数映射在零点附近是局部微分同胚。因此一个等距映射若固定一点且微分在该点为恒等，就在该点的一个邻域恒等。具有这两项性质的点集既闭又开；从已经固定的开集出发，连通性使它等于 \(M\)。作用有效，故该元素只能是单位元。

	于是 \(\chi\) 的微分单射；定义域和值域的 Lie 代数维数都是 \(h\)，故它的像含单位元邻域。紧性使像闭，而 \((S^1)^h\) 连通，像必为整个环面。因此 \(H\cong(S^1)^h\)，特别地 \(H\) 连通，且各 \(\chi_j\) 构成特征格的整数基。当 \(h=0\) 时，同一论证说明稳定子平凡。

	选择 \(H\) 的整数格补，局部模型可写为
	\begin{equation}\label{g4:eq:1204:corner-model}
		T^{n-h}\times\mathbb R^{n-h}_p\times\mathbb C^h_z,
		\qquad
		\omega=\sum_{a=1}^{n-h}\mathrm{d}\theta_a\wedge\mathrm{d}p_a
		+\sum_{j=1}^h\mathrm{d}x_j\wedge\mathrm{d}y_j,
	\end{equation}
	其中 \(H\) 以各坐标的正旋转作用，\(z_j=x_j+\mathrm{i}y_j\)。相对于相应的整数基，矩变量为 \((p,-|z_1|^2/2,\ldots,-|z_h|^2/2)\)，另加 \(\mu(x)\) 的平移。一个矩值确定全部 \(p_a\) 和 \(|z_j|\)，因而局部矩纤维恰为一条轨道。这也给出局部轨道商的象限坐标。

	每条轨道因此在其矩纤维中相对开。\Cref{g4:thm:1202:connected-fibres}说明矩纤维连通，故不能分成两条以上互不相交的开轨道，纤维只能是一条轨道。\(T\) 紧，轨道商是紧 Hausdorff 空间；诱导的连续双射到 \(P\subset\mathfrak t^*\) 因而为同胚。

	\Cref{g4:thm:1202:convexity}已给出 \(P\) 为凸多面体。局部模型中的 \(h\) 个半轴与 \(n-h\) 个自由坐标说明它满维，并在该点有恰好 \(h\) 个相交的边界面。顶点恰好对应 \(h=n\)。此时向内部的原始边向量是切片权的相反数，构成整数基；等价地，各面的原始外法向也构成对偶整数基。这正是 Delzant 条件。一般面上的稳定子由对应边界面的法向圆群生成。
\end{proof}

两个辛环面流形若有同一个矩多面体 \(P\)，上面的结论在每个 \(b\in P\) 给出相同的稳定子、切片权及局部矩锥。一般轨道标准形遂给出 \(b\) 上方轨道附近保持矩映射的 \(T\) 等变辛同构。我们可以把它的定义域取为 \(\mu^{-1}(U)\)，其中 \(U\) 是 \(P\) 中的相对开邻域：若一个轨道邻域包含 \(\mu^{-1}(b)\)，其补集紧，补集的矩像不含 \(b\)，所以缩小 \(U\) 后整条邻近纤维都落在该轨道邻域中。以后取有限个这样的小凸邻域覆盖 \(P\)。

\ProseParagraphBreak

这里的局部同构保存固定的环面 \(T\)。在顶点改用整数格基只是表示同一个局部模型，不能把不同顶点处任意的实线性换基当成全局允许的群变换。

\subsection{自由轨道部分的粘合}
\label{g4:subsec:1204:02}
先在 \(P\) 的内部观察粘合误差。此处作用自由，\(\mu\) 的分量是全局给定的作用变量。在一个小凸集 \(U\subset\operatorname{int}P\) 上，第十一章给出
\[
	\mu^{-1}(U)\cong T\times U,\qquad
	\omega=\sum_{j=1}^n\mathrm{d}\theta_j\wedge\mathrm{d}b_j,
\]
其中 \(b=(b_1,\ldots,b_n)\) 是固定的矩坐标。与 \(T\) 作用交换且保持 \(b\) 的映射只能是
\begin{equation}\label{g4:eq:1204:phase-translation}
	(\theta,b)\longmapsto(\theta+a(b),b),\qquad a:U\to T.
\end{equation}
在 \(U\) 上取 \(a\) 的实提升 \(\widetilde a:U\to\mathfrak t\)。代入辛形式，映射保辛当且仅当
\begin{equation}\label{g4:eq:1204:closed-phase}
	\sum_j\mathrm{d}\widetilde a_j\wedge\mathrm{d}b_j=0,
	\quad\text{即}\quad
	\alpha=\sum_j\widetilde a_j\,\mathrm{d}b_j\text{ 闭}.
\end{equation}
凸性给出 \(\alpha=\mathrm{d}f\)，因此该映射是 \(f\circ\mu\) 的时间 \(1\) Hamilton 流。

\ProseParagraphBreak

这个表达把局部同构之间的误差变成了可相加的函数。记
\begin{equation}\label{g4:eq:1204:hamilton-automorphism}
	\Phi_f(x)=\exp\bigl(\mathrm{d}f_{\mu(x)}\bigr)\cdot x,
\end{equation}
这里用 \(\mathrm{T}_b^*\mathfrak t^*\cong\mathfrak t\) 识别微分，\(\exp:\mathfrak t\to T\) 的核为 \(2\pi\Lambda\)。有 \(\Phi_f\Phi_g=\Phi_{f+g}\)。在连通 \(U\) 上，\(\Phi_f\) 为恒等，当且仅当
\begin{equation}\label{g4:eq:1204:affine-kernel}
	f(b)=c+2\pi\langle b,m\rangle,
	\qquad c\in\mathbb R,\quad m\in\Lambda.
\end{equation}
必要性来自自由性：\(\mathrm{d}f\) 只能取在离散集 \(2\pi\Lambda\) 中，连续性使它恒定；充分性则由指数映射直接得到。

若只在 \(\operatorname{int}P\) 处理这些函数，所得同构还未必在收缩的角圆处光滑。下一小节将证明，光滑的边界模型恰好保证相位及其生成函数也能光滑延至多面体各面。

\subsection{边界轨道的延拓}
\label{g4:subsec:1204:03}
设 \(U\) 为多面体的相对开集，\(f\) 为 \(U\) 上的函数。若 \(f\) 在每一点附近都是环境 \(\mathfrak t^*\) 上某个光滑函数的限制，则称 \(f\) 在 \(U\) 上光滑，记为 \(f\in C^\infty(U)\)。局部象限中的各阶单侧导数也可描述同一个光滑结构。下面先解释从复坐标到象限坐标所需的光滑性事实。

\begin{lemma}\label{g4:lem:1204:radial-division}
	一个在 \(\mathbb C^h\times\mathbb R^k\) 的原点附近光滑、对各复坐标的独立旋转不变的函数，可局部写成
	\[
		A(|z_1|^2,\ldots,|z_h|^2,p),
	\]
	其中 \(A\) 在象限上光滑。若复值光滑函数 \(G_j\) 只在第 \(j\) 个圆作用下具有权 \(1\)，在其他圆作用下不变，则
	\[
		G_j(z,p)=z_j A_j(|z_1|^2,\ldots,|z_h|^2,p)
	\]
	且 \(A_j\) 同样光滑。
\end{lemma}
\begin{proof}
	先处理一个复变量，并允许其余变量作参数。旋转不变函数在实轴的限制 \(g(r,p)\) 是光滑偶函数，其值决定整个函数。对任意 \(N\)，带参数的 Taylor 公式给出
	\[
		g(r,p)=\sum_{j=0}^{N}a_j(p)r^{2j}+r^{2N+2}R_N(r,p),
	\]
	其中余项及任意预先指定的有限阶参数导数局部有界。令 \(s=r^2\)，并对 \(s>0\) 反复使用 \(\partial_s=(2r)^{-1}\partial_r\)。先选足够高阶的 Taylor 展开，就能验证任意有限阶 \(s,p\) 导数都连续延至 \(s=0\)，其边界值为相应 Taylor 系数。故 \(A(s,p)=g(\sqrt s,p)\) 在闭半轴上各阶光滑。

	这也确实给出环境光滑延拓。可取固定截断函数 \(\chi\)，对给定的边界系数逐项构造
	\[
		\sum_{j\ge0}a_j(p)s^j\chi(s/\varepsilon_j).
	\]
	在一个较小紧参数邻域上，逐次选 \(\varepsilon_j\downarrow0\)，使第 \(j\) 项在所有总阶小于 \(j\) 的导数范数中不超过 \(2^{-j}\)。有限低阶项无此限制。该级数在双侧光滑且具有指定的全部边界导数；从 \(A\) 减去它后得到在 \(s=0\) 各阶为零的函数，将差值在 \(s<0\) 延为零即得光滑延拓。这一构造局部适用于参数。

	对权 \(1\) 的 \(G\)，实轴限制满足 \(G(-r,p)=-G(r,p)\)，且 \(G(0,p)=0\)。公式
	\[
		\frac{G(r,p)}r=\int_0^1\partial_rG(tr,p)\,\mathrm{d}t
	\]
	把商光滑延至零；商为偶函数，由前段写成 \(A(r^2,p)\)。旋转等变性再给出 \(G(z,p)=zA(|z|^2,p)\)。对多个复变量，依次应用同一带参数论证；各次取商和平方径向替换均保持其他旋转的不变性，得到所述结论。
\end{proof}

\begin{proposition}\label{g4:prop:1204:boundary-automorphisms}
	设 \(U\subset P\) 是足够小的凸相对开集。任一保持矩映射的 \(T\) 等变辛自同构 \(a:\mu^{-1}(U)\to\mu^{-1}(U)\) 都可写成 \(a=\Phi_f\)，其中 \(f\in C^\infty(U)\)。两个这样的函数给出同一个映射，当且仅当它们相差式\eqref{g4:eq:1204:affine-kernel}中的函数。
\end{proposition}
\begin{proof}
	在内部自由轨道上，相位 \(g(b)\in T\) 唯一确定。考察式\eqref{g4:eq:1204:corner-model}的一张边界图。等变性和保持矩值说明，映射的自由角部分是 \(\theta\mapsto\theta+u(|z|^2,p)\)，复坐标部分满足
	\[
		z_j\longmapsto z_j A_j(|z|^2,p),\qquad |A_j|=1.
	\]
	这里自由角的圆值相位及 \(A_j\) 的光滑性由\cref{g4:lem:1204:radial-division}得到；\(|A_j|=1\) 先在 \(z_j\ne0\) 由模长不变得到，再以连续性延至 \(z_j=0\)。所以内部的 \(g\) 光滑延至全部边界。两张图得到的延拓在稠密内部相同，故在交集上相同。这证明 \(g:U\to T\) 是全局光滑的。

	凸集上的圆值映射有光滑实提升。具体地，对各分量 \(g_j\)，闭一形式 \(-\mathrm{i}g_j^{-1}\mathrm{d}g_j\) 沿从一个内点到 \(b\) 的线段积分，得到 \(\widetilde g_j(b)\)，再调整其在基点的常数，就有 \(e^{\mathrm{i}\widetilde g_j}=g_j\)。线段积分在各面也光滑。内部的保辛条件为式\eqref{g4:eq:1204:closed-phase}，光滑性使它在边界仍成立。再次沿线段积分，得到 \(\sum_j\widetilde g_j\,\mathrm{d}b_j=\mathrm{d}f\)。故 \(a=\Phi_f\) 在内部相同，并由连续性在全部 \(\mu^{-1}(U)\) 相同。

	反过来，\(f\circ\mu\) 在流形上光滑，其 Hamilton 场沿紧轨道切向，故时间 \(1\) 流存在，并等于式\eqref{g4:eq:1204:hamilton-automorphism}。核的刻画在内部已证，函数的连续性将同一仿射表达延至边界。
\end{proof}

边界处单独一条轨道不能唯一确定相位，因为稳定子中的元素对它作用平凡；邻近自由轨道及光滑延拓却确定了同一相位函数。这个区别使我们能够把全部粘合数据放在多面体上处理。

\subsection{等变辛同构与平移自由}
\label{g4:subsec:1204:04}
还需要一条关于凸集覆盖的具体代数事实。它将消去局部 Hamilton 函数在三重交集上的整数仿射误差。

\begin{lemma}\label{g4:lem:1204:convex-cech}
	设 \(P\) 为非空凸多面体，\(\mathcal U=(U_i)\) 为有限相对开凸覆盖，\(K\) 为任意阿贝尔群。以局部常值 \(K\) 系数构成的这份覆盖的 Čech 复形，在正次数的上同调为零。特别地，每个常值 \(K\) 系数的二上圈都是一上链的余边界。
\end{lemma}
\begin{proof}
	每个非空有限交 \(U_{i_0\cdots i_p}\) 都凸。考虑双复形
	\[
		B^{p,q}=\prod_{i_0<\cdots<i_p}C^q(U_{i_0\cdots i_p};K),
	\]
	横向为交替限制，纵向为奇异上边界，总微分取 \(\delta_{\mathrm C}+(-1)^p\delta_{\mathrm s}\)。对固定交集，沿线段向其中一点作锥，给出增广奇异链的收缩；对偶后可见纵向上同调只在零次为 \(K\)。第二卷引理7.2.2的第一象限比较于是将总上同调识别为所述常值 Čech 上同调。

	另一方面，用覆盖小奇异上链增广各横行。对每条小单形 \(\sigma\)，选一个容纳其全部像的指标 \(j(\sigma)\)。在上链的交替指标首位插入 \(j(\sigma)\)，逐项求值即给出横向收缩；交替和展开后满足 \(\delta h+h\delta=1\)。因此同一比较引理把总上同调识别为小奇异上链的上同调。这里无须令小单形的指标选择随 \(\sigma\) 连续变化。

	第二卷定理2.3.8说明小链包含在整数同调上是同构；小链和全部奇异链都为自由整数链复形，第二卷定理4.2.3的自然上同调万有系数列因而把此同构转成任意 \(K\) 系数的上同调同构。最后，\(P\) 自身也可沿线段收缩到一点，锥算子使其正次上同调为零。结合两种比较，即得结论。这个证明只用了凸性和相对开覆盖，所以也适用于 \(P\) 的边界与角点。
\end{proof}

\begin{theorem}[Delzant 分类的唯一性]\label{g4:thm:1204:uniqueness}
	设 \((M_i^{2n},\omega_i,T,\mu_i)\)、\(i=1,2\)，为两个辛环面流形，环面及其整数格已固定。若 \(\mu_1(M_1)=\mu_2(M_2)=P\)，则存在 \(T\) 等变辛同胚 \(\Psi:M_1\to M_2\)，使 \(\mu_2\Psi=\mu_1\)。
\end{theorem}
\begin{proof}
	由\cref{g4:prop:1204:toric-orbits}及其后的局部比较，取有限相对开凸覆盖 \((U_i)\) 和局部同构
	\[
		\psi_i:\mu_1^{-1}(U_i)\longrightarrow\mu_2^{-1}(U_i).
	\]
	这些集合可取为 \(P\) 与小 Euclid 球的交，因而全部非空有限交仍凸。对两重交集，定义源流形上的转移自同构 \(a_{ij}=\psi_i^{-1}\psi_j\)。\Cref{g4:prop:1204:boundary-automorphisms}在局部给出光滑相位；其在内部的唯一性使这些相位在整个交集相容。交集凸，故同一命题的积分步骤在整个交集给出 \(f_{ij}\in C^\infty(U_i\cap U_j)\)，满足
	\[
		a_{ij}=\Phi_{f_{ij}},\qquad f_{ji}=-f_{ij}.
	\]
	转移映射满足 \(a_{ij}a_{jk}=a_{ik}\)，所有这类自同构交换，因此
	\[
		f_{ij}+f_{jk}-f_{ik}
		=c_{ijk}+2\pi\langle b,m_{ijk}\rangle,
		\qquad (c_{ijk},m_{ijk})\in\mathbb R\oplus\Lambda.
	\]
	每个交集连通，右端的系数为常量。它们是一个 \(K=\mathbb R\oplus\Lambda\) 系数的二上圈，因为左端是 \(f\) 的 Čech 余边界，而 \(\delta^2f=0\)。仿射函数在非空交集中的表达唯一：每个这样的相对开交集都含有 \(P\) 的内点，不能使一个非零仿射函数恒为零。

	由\cref{g4:lem:1204:convex-cech}，可找到常整数仿射的一上链，其余边界恰等于上述误差。从 \(f_{ij}\) 减去该一上链不会改变 \(\Phi_{f_{ij}}\)，却使
	\begin{equation}\label{g4:eq:1204:exact-function-cocycle}
		f_{ij}+f_{jk}=f_{ik}.
	\end{equation}
	这一步同时处理实常数和整数周期，因而后续可以对真正的函数余圈使用单位分解。

	在环境 Euclid 空间取从属于这些小邻域的光滑非负截断，并在 \(P\) 附近以它们的正和归一化，得到 \(P\) 上的光滑单位分解 \((\rho_k)\)，且 \(\operatorname{supp}\rho_k\subset U_k\)。令
	\[
		h_i=\sum_k\rho_k f_{ik}\quad\text{于 }U_i.
	\]
	每项在相应交集定义，再借助支集条件补零；有限求和保证光滑。式\eqref{g4:eq:1204:exact-function-cocycle}给出
	\[
		h_i-h_j=\sum_k\rho_k(f_{ik}-f_{jk})=f_{ij}.
	\]
	所以修正后的局部同构 \(\widetilde\psi_i=\psi_i\Phi_{h_i}\) 满足
	\[
		\widetilde\psi_i^{-1}\widetilde\psi_j
		=\Phi_{-h_i}\Phi_{f_{ij}}\Phi_{h_j}=\operatorname{id}.
	\]
	它们拼成全局 \(\Psi\)。每张图中的映射光滑、保辛、等变并保持矩映射，这些性质在拼接后保持。每条目标纤维上，它由任一覆盖该底点的局部同构给出双射，所以 \(\Psi\) 全局双射；局部逆也相容，故为所述等变辛同胚。
\end{proof}

结合前节的构造，得到整个分类：固定 \(T\)、指定矩映射的紧连通辛环面流形，按保持矩映射的 \(T\) 等变辛同胚分类，与 \(\mathfrak t^*\) 中的 Delzant 多面体一一对应。分类定理的标准陈述可参见 Cannas da Silva 的定理6.8；前节的约化与本节的粘合分别完成存在性和唯一性。\cite[第6.4节，第169--174页]{Cannas2006Symplectic}

\begin{proposition}\label{g4:prop:1204:translation-lattice-freedom}
	在上述分类中，若不指定矩映射的加法常数，则多面体按平移识别。若还允许以环面自同构重标作用，则多面体进一步按保持整数格的线性变换识别，即按整数仿射变换分类。
\end{proposition}
\begin{proof}
	同一个连通辛流形上，同一作用的两个矩映射对每个 \(\xi\in\mathfrak t\) 的分量有相同微分。因此它们的差是常向量 \(c\in\mathfrak t^*\)，矩像相差平移。

	若环面自同构 \(\kappa:T\to T\) 的微分为 \(A\)，则 \(A\Lambda=\Lambda\)，相对于整数基有 \(A\in\operatorname{GL}(n,\mathbb Z)\)。若 \(\Psi(t\cdot x)=\kappa(t)\cdot\Psi(x)\) 且 \(\Psi\) 保辛，则
	\[
		\mathrm{d}\langle\mu_2\Psi,A\xi\rangle
		=\Psi^*\iota_{(A\xi)_{M_2}}\omega_2
		=\iota_{\xi_{M_1}}\omega_1
		=\mathrm{d}\langle\mu_1,\xi\rangle.
	\]
	所以 \(A^{\mathsf T}\mu_2\Psi=\mu_1+c\)。这给出对偶空间中的整数线性变换及平移；反过来，重标群作用及平移矩映射后应用唯一性定理，即得相应同构。
\end{proof}

多面体的尺度并未被上述等价关系消去。例如一维区间长度为 \(\ell\) 时，对应球面的辛面积为 \(2\pi\ell\)，不同长度不能由辛同胚识别。分类保存的是格、各面位置及其相对尺度；只保留多面体的组合类型，会丢失这些辛几何数据。
% 纲要条目排为编号小节标题，并进入目录与书签；说明段仍保留于纲要注释中。
